\documentclass[11pt]{article}

\usepackage[margin=1in]{geometry}
\usepackage[T1]{fontenc}
\usepackage{amsmath, amssymb}
\usepackage{booktabs}
\usepackage{enumitem}
\usepackage{titlesec}
\usepackage{fancyhdr}
\usepackage{float}
\usepackage{xcolor}

\usepackage[hidelinks]{hyperref}

\titlespacing*{\section}{0pt}{1.2em}{0.5em}
\titlespacing*{\subsection}{0pt}{0.8em}{0.4em}
\titleformat{\section}{\large\bfseries}{\thesection.}{0.5em}{}
\titleformat{\subsection}{\normalsize\bfseries}{\thesubsection}{0.5em}{}

\pagestyle{fancy}
\fancyhf{}
\lhead{\small CS 388: Natural Language Processing}
\rhead{\small Homework 2}
\cfoot{\small\thepage}
\renewcommand{\headrulewidth}{0.4pt}

\setlist[enumerate]{itemsep=2pt, topsep=3pt}
\setlist[itemize]{itemsep=2pt, topsep=3pt}

\newcommand{\wrd}[1]{\texttt{#1}}
\newcommand{\emb}[1]{\ensuremath{e(#1)}}
\newcommand{\sm}{\ensuremath{\mathrm{softmax}}}
\newcommand{\Wq}{\ensuremath{W_Q}}
\newcommand{\Wk}{\ensuremath{W_K}}
\newcommand{\Wv}{\ensuremath{W_V}}
\newcommand{\pred}{\ensuremath{\hat{y}}}

\definecolor{solblue}{RGB}{0,60,180}
\newcommand{\sol}[1]{\par\medskip{\color{solblue}\noindent\textbf{Solution.} #1}\par\medskip}
\newcommand{\B}[1]{{\color{solblue}#1}}

\begin{document}

\begin{center}
  {\Large\bfseries Homework 2: Self-Attention and Transformers}\\[2pt]
  {\normalsize CS 388: Natural Language Processing \quad\textbullet\quad Fall 2026}\\[2pt]
\end{center}
\vspace{0.5em}
\hrule
\vspace{1em}

\noindent
\textbf{Instructions.} This assignment can be done entirely by hand -- no coding required (Question 4 asks you to read a short PyTorch file and describe a fix, but you never have to run anything); however, please transcribe your answers/work into the provided \LaTeX file and submit the PDF, rather than a handwritten response.
\textbf{Please use the {\tt{\textbackslash B\{\}}} command to wrap your solutions, making them blue (this will help us when grading).}
Show your work. When you are asked to construct a matrix, write out every entry, and give each entry as a
number. Submit a single PDF to Gradescope.

\section*{Setup and Notation}

Our alphabet is $\Sigma = \{\wrd{A}, \wrd{B}, \wrd{C}\}$, embedded as one-hot vectors in $\mathbb{R}^3$:
\[
\emb{\wrd{A}} = [1,0,0], \qquad \emb{\wrd{B}} = [0,1,0], \qquad \emb{\wrd{C}} = [0,0,1].
\]
All vectors are row vectors. A sequence is a string $x = x_1 x_2 \cdots x_n$ with each $x_i \in \Sigma$; we write
$X \in \mathbb{R}^{n \times 3}$ for the matrix whose $i$-th row is $\emb{x_i}$. Throughout this assignment
every sequence has length at most $L = 5$.

\medskip\noindent
\textbf{Model.}
We will use a single layer of scaled dot-product self-attention, with parameter matrices
$\Wq$, $\Wk$, and $\Wv$, each in $\mathbb{R}^{3\times 3}$. For each position $i$ we form a query, a key,
and a value,
\[
q_i = \emb{x_i}\Wq, \qquad k_i = \emb{x_i}\Wk, \qquad v_i = \emb{x_i}\Wv ,
\]
and the attention output at position $i$ is $\sm\!\left(q_i K^\top\right) V$. Real implementations divide the
scores by $\sqrt{d}$ to stop them from growing with the embedding dimension; here $d$ is tiny and we are
choosing the weights by hand, so we drop that constant --- it would only rescale $\Wq$.

We only ever care about the output at the \emph{final} position, which is the position that has read the whole
string. So write $q = q_n = \emb{x_n}\Wq$, let
\[
s_i = q\, k_i^{\top}, \qquad
\alpha_i = \frac{\exp(s_i)}{\sum_{j=1}^{n}\exp(s_j)}, \qquad
o \;=\; \sum_{i=1}^{n} \alpha_i v_i \;\in\; \mathbb{R}^3 ,
\]
so $s_i$ is the score of position $i$, $\alpha$ is the attention distribution, and $o$ is the output vector.
Note that position $n$ attends to itself: the sum runs over all of $1,\dots,n$, including $i = n$.

\medskip\noindent
\textbf{Inference.} We read the three coordinates of $o$ off as scores for the three symbols, and the model's
prediction is
\[
\pred(x) \;=\; \operatorname*{arg\,max}_{\ell \in \{\wrd{A},\wrd{B},\wrd{C}\}} \; o_\ell .
\]
Initially, we will not have any output projection or MLP (until Question 2, where we add an output
projection). A tie for the argmax counts as a wrong answer.
A rule is \emph{captured} by a choice of parameters if $\pred(x)$ is correct for
\emph{every} sequence $x$ of length $n \le L$ to which the rule applies.

\medskip\noindent
\textbf{Simplification.} In Questions 1 and 2 we will fix $\Wk = \Wv = I_3$, so that $k_i = v_i = \emb{x_i}$,
leaving $\Wq$ as the only parameter for you to design. Under this choice, if $x_n$ is the $a$-th symbol of the alphabet and $x_i$ is the
$b$-th, then
\[
s_i \;=\; \emb{x_n}\,\Wq\,\emb{x_i}^{\top} \;=\; (\Wq)_{ab} ,
\]
i.e.\ $\Wq$ is nothing more than a $3\times 3$ lookup table of scores: the entry \emph{is} the score, indexed by (final symbol, attended
symbol). Feel free to use this fact without re-deriving it.

\medskip\noindent
\textbf{Reading off the output.} Because $\Wv = I_3$, the values are the embeddings themselves, so
$o = \sum_i \alpha_i \emb{x_i}$, and the $\ell$-th coordinate of $o$ is just the total attention mass placed on
the positions holding $\ell$:
\[
o_\ell \;=\; \sum_{i \,:\, x_i = \ell} \alpha_i .
\]
In particular the coordinates of $o$ are nonnegative and sum to $1$, so $o$ is a probability distribution over
$\Sigma$. You may use both of these facts without proof; you will need them in every question.

\section{Question 1 (25 points)}

Consider sequences over $\{\wrd{A}, \wrd{B}\}$ only (no \wrd{C} appears), and the following pair of decision
rules:
\begin{center}
\begin{tabular}{ll}
\toprule
R1 & if the sequence is all \wrd{A}'s, predict \wrd{A}\\
R2 & if the sequence contains at least one \wrd{B}, predict \wrd{B}\\
\bottomrule
\end{tabular}
\end{center}
Together these cover every sequence over $\{\wrd{A},\wrd{B}\}$. Note that R2 has no count threshold: a single
\wrd{B} among a thousand \wrd{A}'s still means \wrd{B}.

\begin{enumerate}[label=(\alph*)]

  \item \textbf{(15 pts)} Write down a matrix $\Wq$ that captures R1 and R2, giving all nine entries.

  \item \textbf{(10 pts)} Verify your answer on the three sequences \wrd{AAA}, \wrd{AAB}, and \wrd{ABA} by
  computing $s_i$, $\alpha_i$, and $o$ for each.



\end{enumerate}

\section{Question 2 (40 points)}

Now let all three symbols appear, and consider the following rule set, which extends the previous one:
\begin{center}
\begin{tabular}{lccl}
\toprule
& contains \wrd{B}? & contains \wrd{C}? & prediction\\
\midrule
R1 & no  & no  & \wrd{A}\\
R2 & yes & no  & \wrd{B}\\
R3 & no  & yes & \wrd{A}\\
R4 & yes & yes & \wrd{C}\\
\bottomrule
\end{tabular}
\end{center}
In other words, we predict \wrd{B} if a sequence contains a \wrd{B} and no \wrd{C}; we predict \wrd{C} if and
only if it contains both a \wrd{B} and a \wrd{C}; and otherwise we predict \wrd{A}.
The rules still cover every sequence, and we still fix $\Wk = \Wv = I_3$.

\begin{enumerate}[label=(\alph*)]

  \item \textbf{(12 pts)} Show that no single $\Wq$ captures R1--R4. \emph{Hint:} both \wrd{CA} and \wrd{BCA}
  end in the same symbol, so they use the same row of $\Wq$, and they contain the same number of \wrd{A}'s and
  the same number of \wrd{C}'s. You should be able to write out two incompatible inequalities that these rules require. 


  \item \textbf{(6 pts)} Explain in two or three sentences \emph{why} this happens, in terms of how the scores
  $s_i$ are computed. Your explanation should make clear what property of the rule set a single head cannot
  represent, and should not depend on the particular counterexample from (a).

  \item \textbf{(16 pts)} Still assuming $\Wk = \Wv = I_3$, show that multi-head attention with two heads ---
  that is, two separate query matrices --- can capture the rules. Heads $1$ and $2$ should produce outputs $o^{(1)}, o^{(2)} \in
  \mathbb{R}^3$. These are concatenated into $[o^{(1)}; o^{(2)}] \in \mathbb{R}^6$ and passed through a linear output layer
  \[
  z \;=\; [o^{(1)}; o^{(2)}]\, W_O, \qquad W_O \in \mathbb{R}^{6\times 3}\;,
  \]
  with the prediction now $\pred(x) = \arg\max_\ell z_\ell$. Design $\Wq^{(1)}$, $\Wq^{(2)}$, and $W_O$ so that some
  coordinate of $o^{(1)}$ behaves as an (approximate) indicator of ``$x$ contains \wrd{B}'' and some coordinate
  of $o^{(2)}$ as an indicator of ``$x$ contains \wrd{C}''. Give the matrices and identify the two coordinates.
  Remember that your construction has to be correct for every sequence of length at most $L = 5$, not just for
  short ones.

  \item \textbf{(6 pts)} Why can the rules above be captured by a model with no non-linearity? 
  

\end{enumerate}

\section{Question 3 (20 points)}

For this part, restrict the alphabet to $\{\wrd{A}, \wrd{C}\}$ and consider a single rule:
\begin{center}
\begin{tabular}{ll}
\toprule
R5 & if the first symbol is \wrd{C}, predict \wrd{C}; otherwise predict \wrd{A}\\
\bottomrule
\end{tabular}
\end{center}

\begin{enumerate}[label=(\alph*)]

  \item \textbf{(5 pts)} Can single-head or multi-head attention, as we have defined it so far, capture this
  rule? Why or why not? (No proof needed here, 1--2 sentences of explanation is fine.)

  \item \textbf{(15 pts)} Fix the model by hand with an integer positional encoding. Embed position $i$ by
  appending its index, so that token $i$ is represented by the $4$-dimensional vector
  \[
  \tilde{e}_i \;=\; [\,\emb{x_i},\; i\,] \;\in\; \mathbb{R}^4 ,
  \]
  take $\Wq, \Wk, \Wv \in \mathbb{R}^{4\times 4}$, and let the prediction be the argmax over the
  first three coordinates of $o$ as before. Give $\Wq, \Wk, \Wv$ explicitly (all sixteen entries of each)
  such that R5 is captured for all sequences of length at most $L = 5$.




\end{enumerate}

\section{Question 4 (15 points)}

The file \texttt{mhsa.py} accompanying this assignment implements a single layer of multi-head self-attention in
PyTorch. It is meant to do exactly what you did by hand in Question 2: run each head independently, recombine
the per-head outputs, and project the result back to the model dimension with $W_O$. The file contains
\textbf{exactly one} mistake --- the shapes, the scaling, the softmax, and the four projections are all
correct. You do not need to run anything.

\begin{enumerate}[label=(\alph*)]

  \item \textbf{(5 pts)} Find the mistake. Quote the offending line, and say in one sentence what it does
  instead of what it was supposed to do.

  \item \textbf{(5 pts)} What happens if the mistake is not fixed? Describe what you would see if you ran the
  file as given, and also what the layer would actually be computing if someone patched over that symptom
  rather than fixing the real problem. Connect the second half of your answer to what you showed in
  Question 2(a).

  \item \textbf{(5 pts)} Give a one-line fix: the single line of code that should replace the offending line.

\end{enumerate}

\vspace{1em}
\hrule
\vspace{0.5em}
\noindent
\small
\textbf{What to turn in.} A single PDF containing every matrix you were asked to construct, written out
entry by entry; the worked verifications requested in Question 1(b); the impossibility argument in
Question 2(a); and your written answers to the discussion questions. Do not submit code -- write the one-line
fix from Question 4(c) directly into your PDF.

\end{document}
